\scanpage{43}\mspage{$\underline{\overline{\mathrm{XV}}}$}
Choisissons un voisinage $V$ de $\Gamma_{x'y'}$ tel que : $\mathrm{I})\ V\cap\Gamma_{ss'}=\varnothing$ ;
% Paired vertical strokes at the two edges of the upper sheet.
\reversemarginpar\margin{$\vert$}\normalmarginpar\margin{$\vert$}%
$\overline{\underline{\mathrm{II}}})\ V\cap\mathrm{int}\,\Lambda^{2}$, $V\cap\mathrm{ext}\,\Lambda^{2}$ sont des ouverts connexes. (Il \unsure{\.{y}|\'{y}|y} en a toujours car si $A$ et $B$ sont deux parties compactes de $\mathbb{R}^{2}$ il \unsure{exìste|existe} un couple de voisinages $A\subset V_A$, $B\subset V_B$ tel que $V_A\cap V_B=\varnothing$, le lemme I \deleted{\unsure{fai|fait}} montre qu'on peut trouver $V\inserted{\subset V_A}$ satisfaisant en plus $\overline{\underline{\mathrm{II}}})$ si $A=\Gamma_{\deleted{\unsure{s}}x'\deleted{\unsure{s}}y'}$ et $B=\Gamma_{ss'}$)

Comme $\Lambda_{xx'}\cap V\inserted{-\{x'\}}\ne\varnothing$ \deleted{\illegible{1}} (resp. $\Lambda_{y'y}\cap V-\{y'\}\ne\varnothing$) on a
\[
\begin{aligned}
&\Lambda_{xx'}\cap V-\underbrace{\Lambda_{xx'}\cap\Lambda^{2}}_{x'}\subset V\cap\mathrm{int}\,\Lambda^{2}\ \cup\ V\cap\mathrm{ext}\,\Lambda^{2}\quad(\text{resp.}\\
&\Lambda_{y'y}\cap V-\underbrace{\Lambda_{y'y}\cap\Lambda^{2}}_{y'}\subset V\cap\mathrm{int}\,\Lambda^{2}\ \cup\ V\cap\mathrm{ext}\,\Lambda^{2})\quad\text{mais alors}\\
&\Lambda_{xx'}\cap V-\{x'\}\subset V\cap\mathrm{ext}\,\Lambda^{2}\quad(\text{resp. }\Lambda_{y'y}\cap V-\{y'\}\subset V\cap\mathrm{ext}\,\Lambda^{2})\\
&\text{car}\quad \Lambda_{xx'}-\{x'\}\subset\mathrm{ext}\,\Lambda^{2}\quad(\text{resp. }\Lambda_{y'y}-\{y'\}\subset\mathrm{ext}\,\Lambda^{2}).
\end{aligned}
\]
Soit donc $x'',\deleted{y''\in}\ \in(\Lambda_{xx'}-\{x'\})\cap V$ (resp. $y''\in(\Lambda_{y'y}\unsure{\leftarrow|-}\{y'\})\cap V$).
Il est clair qu'on peut rélier $x''$ à $y''$ par un arc de polygone $\inserted{\Lambda_{x''y''}}\subset V$ avec $\Lambda_{x''y''}\cap\Gamma^{\prime}_{ss'}=\varnothing$ (cela résulte \unsure{dee|deu|du} fait que $V\cap\Gamma'_{ss'}=\varnothing$) alors l'arc de polygone $\inserted{\overline{\Lambda}_{xy}}$ obtenu
\reversemarginpar\margin{$\vert$}\normalmarginpar\margin{$\vert$}%
en récollant $\Lambda_{xx''}$ avec $\Lambda_{x''y''}$, $\Lambda_{x''y''}$ avec $\Lambda_{y''y}$, selon $\{x''\}$ et $\{y''\}$ respectivement, satisfait $\overline{\Lambda}_{x\dot y}\subset U=\unsure{\mathrm{Int}|\mathrm{int}}\,\Lambda^{1}-R_{\Lambda^{2}}\cap\mathrm{int}\,\Lambda^{1}$.

Prouvons que $V$ est connexe. Soient $x,y\in U=\mathrm{ext}\,\Lambda^{1}\cup\mathrm{int}\,\Lambda^{2}\cup{]s,s'[}$

Si $x,y\in\mathrm{ext}\,\Lambda^{1}$ on rélie $\dot{x}$ à $y$ par un arc de polygone $\Lambda_{xy}\subset\mathrm{ext}\,\Lambda^{1}$.

Si $x,y\in\mathrm{int}\,\Lambda^{2}\cup{]s,s'[}$ on peut aussi rélier $x$ à $y$ par un arc de polygone $\Lambda_{xy}\subset\mathrm{ext}\,\Lambda^{1}$.

Donc il reste le cas ou $x\in\mathrm{ext}\,\Lambda^{1}$ et $y\in\mathrm{int}\,\Lambda^{2}\cup{]s,s'[}$. On choisit $z\in{]s,s'[}$ on rélie $x$ a $z$ par $\Lambda_{xz}$ tel que $\Lambda_{xz}-\{z\}\subset\mathrm{ext}\,\Lambda^{1}$ et $y$ à $z$ par $\Lambda_{zy}$ tel que $\Lambda_{zy}\deleted{\unsure{-\{z\}}}\subset\mathrm{int}\,\Lambda^{2}\cup{]s,s'[}$. Il est \deleted{\unsure{q}}clair \unsure{quén|qu'en} récollant $\Lambda_{xz}$ avec $\Lambda_{zz}$ selon $\{z\}$ on obtient $\Lambda_{xy}\subset V$.

On a vu que $U$ et $V$ sont des connexes il sont necessairement ouverts car il sont les composantes connexes de $\mathbb{R}^{2}-\Gamma$, et $\Gamma$ est compact donc fermé dans $\mathbb{R}^{2}$\unsure{)|,|;}

\hfill c.q.s.d.
% archive sheet number: 42
